\documentclass[11pt]{article}
\usepackage[margin=1in]{geometry}
\usepackage{hyperref}
\usepackage{enumitem}
\title{AI for Oceans and Robotics}
\author{Piter Garcia}
\date{\today}
\begin{document}
\maketitle

\section*{Grade and Class Match}
\textbf{Grade:} Grade 2\\
\textbf{Likely blocks:} Day 1 · 10:45-11:35 · Callon; Day 1 · 2:15-3:05 · Henchen; Day 4 · 10:45-11:35 · Barthelman

\section*{Lesson Focus}
Machine-learning ideas, sorting, and early robotics connections for Grade 2.

\section*{Learning Targets}
\begin{itemize}[leftmargin=*]
\item Students describe AI or machine-learning ideas in simple language.
\item Students complete a guided digital classification or coding task.
\item Students connect the activity to a real-world example.
\end{itemize}

\section*{Materials}
\begin{itemize}[leftmargin=*]
\item Student devices
\item AI or coding platform
\item Projected model
\end{itemize}

\section*{Suggested Lesson Flow}
\begin{enumerate}[leftmargin=*]
\item Open with the exact device, tool, or routine students need in order to start successfully.
\item Model one example task before students begin independent or partner work.
\item Give students time to build, code, design, or document their work while circulating for support.
\item Close with a short share-out, reflection, or debugging discussion.
\end{enumerate}

\section*{Source Notes}
This lesson packet is enriched with transcript-derived lesson notes, the school schedule roster, and official starting-point resources. See the paired Markdown guide for clickable links.

\bibliographystyle{plain}
\bibliography{ai-for-oceans-and-robotics}
\end{document}
